\subsection{Reduktion kognitiver Last durch Struktur}
\label{subsec:02-03-02_cognitive-load}

\Gls{cognitiveload} entsteht, wenn Nutzende Informationen aktiv verarbeiten, Entscheidungen treffen oder Zusammenhänge erschließen müssen und dabei die begrenzte Kapazität des Arbeitsgedächtnisses beansprucht wird \cite{105:Sweller1988,414:Johnson2021}. Eine ungünstig gestaltete \Gls{ia} kann diese Belastung erhöhen, etwa durch tief verschachtelte Navigationsstrukturen, uneinheitliche Benennungen oder redundante Inhalte, die extrinsische kognitive Belastung erzeugen. Darunter wird jene zusätzliche mentale Anstrengung verstanden, die nicht aus der eigentlichen Arbeitsaufgabe entsteht, sondern durch die Struktur und Darstellung des Systems verursacht wird und damit Ressourcen von den eigentlichen Arbeitsaufgaben abzieht \cite{105:Sweller1988,414:Johnson2021}.

Eine strukturierte und klar gegliederte \Gls{ia} reduziert kognitive Last, indem sie Orientierung bietet, Auswahlmöglichkeiten begrenzt und relevante Inhalte sichtbar priorisiert \cite{414:Johnson2021,415:Krug2014}. Prinzipien wie sinnvolle Gruppierung, Hierarchisierung und progressive Offenlegung unterstützen Nutzende dabei, Informationen schrittweise zu erfassen und Aufgaben effizient zu bearbeiten, ohne von irrelevanten Details abgelenkt zu werden \cite{415:Krug2014}. Dadurch verbessert sich nicht nur die Nutzbarkeit, sondern auch die wahrgenommene Nützlichkeit des Systems, da Aufgaben mit geringerem kognitivem Aufwand und höherer Erfolgswahrscheinlichkeit erledigt werden können \cite{102:Davis1989,300:FraunhoferIESE2018}.
